\section{Scheme representation and the exactness gate}\label{sec:repr}

\paragraph{Terms and the MFW1 format.}
A term $(u,v,w)$ is stored as three bit masks: bit $i\,m+j$ of $u$ is $A[i][j]$, bit $j\,p+k$ of $v$ is $B[j][k]$ and bit $i\,p+k$ of $w$ is
$C[i][k]$ (the output layout is row-major and \emph{not} transposed; the public tables that store the third factor column-major are
transposed on import, and a spec asserts that a transposed factor fails the gate).
A scheme is serialized as text (``MFW1''): a header \texttt{MFW1 n m p rank}, then one line \texttt{a b c} per term with the
three masks as lower-case hexadecimal without leading zeros, terms sorted in increasing order of $(a,b,c)$ and each line
terminated by a newline. The parser rejects any input that does not re-serialize to itself and any duplicated term.
The SHA-256 of the canonical bytes identifies a tensor everywhere in this paper; bundle file names embed its first eight hex digits.
A successful parse validates the layout only. The tensor identity is checked separately, by a gate.

\paragraph{Wide factors.}
In the walker each factor is a vector of $\lceil\max(nm,mp,np)/32\rceil\le 32$ little-endian 32-bit limbs, one limb per 64-bit word, so
factors of up to 1{,}024 bits, formats with all sides at least 2 and $nm,mp,np\le1024$, and seed ranks up to 16{,}320 (files of up to 16{,}384 terms) are supported; the same
layout is used for every shape, including those whose factors fit a machine word.
(The project's earlier production fleet uses a narrower one-word representation restricted to $2\le n\le7$; it is not used in this paper.)

\paragraph{The gate.}
\begin{definition}[exact]
A list of triples of nonzero masks is \emph{exact for $\MM nmp$} if $\sum_t u_t\tensor v_t\tensor w_t = T_{n,m,p}$ over $\F$;
coordinatewise, for every $\alpha<nm$ and $\beta<mp$,
\[
  \bigoplus_{t:\;\alpha\in\supp u_t,\;\beta\in\supp v_t} w_t \;=\;
  \begin{cases} 2^{(\alpha\,\mathrm{div}\,m)\,p+(\beta\bmod p)} & \text{if } \alpha\bmod m=\beta\,\mathrm{div}\,p,\\ 0 & \text{otherwise.}\end{cases}
\]
\end{definition}
The rank claimed for a list is its number of terms; the gate does not require the terms to be irredundant, which is harmless for an upper bound.
The sparse-table gate costs $\sum_t|u_t|\,|v_t|$ word operations on masks of at most $np$ bits plus $nmp$ for the target.
Three implementations exist and disagree about nothing on any tensor we tested: a Python function that accumulates each $w_t$ into a
sparse table indexed by $(\alpha,\beta)$, XORs the target into the same table and asserts that it is empty; a Ruby program with its own
parser and its own formulation, which builds the dense $nm\times mp$ table and compares every cell with the closed form above; and the walker's
own packed check in Tungsten. \Cref{tab:verifiers} summarizes what was checked by which.

\begin{table}[t]
\centering\small
\begin{tabular}{@{}p{3.6cm}p{2.1cm}p{7.7cm}@{}}
\toprule
verifier & language & role in the campaign\\
\midrule
walker's packed gate & Tungsten & seed on entry and best scheme on exit of every walk; canonicalization must preserve the rank\\
sparse-table gate & Python & every tensor admitted to the scoreboard (verdicts cached by content hash), every row of every bundle\\
dense-table gate & Ruby & run after the fact over the 691 record tensors held in the archive, all 661 chain heads and all 1{,}089 tensors behind the records (985 distinct bundle tensors, 102 committed certificates, 2 fleet checkpoints): all exact; also inside the leaf check of the composition library. The campaign's walk and projection drivers did not call it\\
\bottomrule
\end{tabular}
\caption{The exactness verifiers. The Ruby program does not invoke any code of the searcher, shares no parser with it, and formulates the
target tensor differently; it shares the bit-layout convention by design, so a convention error common to both would go unnoticed.}
\label{tab:verifiers}
\end{table}

\paragraph{Negative controls.}
On a rank-561 scheme for $\MM3{14}{18}$ both the Python and the Ruby gate accept the original and reject the same scheme with one bit flipped,
with one term dropped, and with the third factor transposed into the public column-major layout. A scheme with one term added twice is
accepted by both (it is still a valid, redundant decomposition of two more terms); such files are rejected earlier by the canonical parser.
The Python gate relies on assertions, so the verifiers that use it refuse to run under \texttt{python -O}.

\paragraph{Cost of checking.}
On the loaded development machine the Python gate took $<0.01$\,s for a $9\times9\times9$ scheme of rank 482, 0.04\,s for
$16\times16\times16$ at rank 2209, and 2.1\,s for a $28\times32\times32$ scheme of rank 14{,}271; the Ruby gate took 0.2\,s, 0.1\,s and 1.2\,s
respectively (the first is dominated by start-up). Checking is never the bottleneck.
